\documentclass[aps,prd,reprint,nofootinbib,longbibliography]{revtex4-2}

\usepackage{amsmath,amssymb,amsthm,bm}
\usepackage{booktabs}
\usepackage{graphicx}
\usepackage{hyperref}
\usepackage{microtype}
\usepackage{pgfplots}
\pgfplotsset{compat=1.15,
  paperaxis/.style={
    width=\columnwidth, height=0.58\columnwidth,
    grid=major, grid style={gray!20},
    tick label style={font=\footnotesize},
    label style={font=\footnotesize},
    legend style={font=\footnotesize, draw=none, fill=none},
    line width=0.9pt},
  discard if not/.style 2 args={
    x filter/.append code={
      \edef\tempa{\thisrow{#1}}
      \edef\tempb{#2}
      \ifx\tempa\tempb\else\def\pgfmathresult{nan}\fi
    }
  }
}

% Figure data files, generated by the committed data script from the
% committed evidence records.
%
% figdata/pf4_cells.dat   columns: run split P E d xg neglogf
%   one row per usable cell of the two governed runs; run is 1 or 2,
%   split is 0 for training and 1 for held-out, xg = (d/E)^2,
%   neglogf = -log(fraction).
% figdata/pf4_gline.dat   columns: xg pred
%   the sealed run-002 G-model line pred = alpha + beta*xg.
% figdata/pf4_manifold.dat columns: P E d regime
%   the eighteen probe points of the pilot record, P written exactly as
%   0.5, 0.7, 0.9; regime is 1 for a deterministically reversing probe
%   and 0 otherwise.

\begin{document}

\title{A Sealed Negative for Schwinger Scaling in a Resolved-Field
Hidden-Dynamics Family}

\author{Andrew H. Bond}
\email{andrew.bond@sjsu.edu}
\thanks{ORCID 0009-0003-2599-6158}
\affiliation{San Jos\'e State University, San Jos\'e, California, USA}
\date{August 4, 2026}

\begin{abstract}
A local hidden dynamics observed through a fold in its projected time can
produce apparent pair events.  Whether it can also produce the field
dependence of the Schwinger exponential, without that law being inserted
by hand, is the sharpest question a campaign of this kind can ask.  This
paper reports the answer for one declared family, under seal, and the
answer is no.  In a resolved Sauter-slab energy-transfer model the
suppression of time reversal is a Gaussian measure tail in a
probe-measured effective gap, and a Schwinger-shaped alternative fails to
compete on held-out data.  The paper reports the full chain that produced
this verdict.  A four-iteration pilot turned each of its failures into a
design constraint.  A first sealed run returned the verdict neither-axis
because its bars, set in binomial units, tested a precision claim nobody
had made, and that verdict is preserved rather than repaired.  An
owner-imposed correction rule, declared in writing before any second
attempt existed, required the next preregistration either to model the
two-variable structure of the exponent explicitly or to set its bars
against model error.  The second sealed run took the second branch and
passed both bars with room.  The effective-gap model predicted gaps it
had never seen better than it fit the gaps it trained on, and the
Schwinger-shaped model was more than three times worse on the same
held-out cells.  The result licenses no statement about other families,
nothing about physical pair production, and nothing about quantum
electrodynamics.  The question it sharpens is whether any constructible
hidden-dynamics family has an exponent that is not a measure tail.
\end{abstract}

\maketitle

\section{Introduction}
\label{sec:intro}

An earlier paper of this campaign proved that a smooth trajectory on a
hidden manifold, observed through a projection into spacetime, produces an
apparent particle-antiparticle pair whenever the observer's time functional
has a nondegenerate critical point along the trajectory \cite{ZenodoEntropy}.
That statement is a theorem about kinematics and it is cheap.  The
expensive question is dynamical.  Can a specified local hidden dynamics,
with no target law inserted, produce reversal statistics with the field
dependence of the Schwinger pair-production exponential
\cite{Sauter1931, HeisenbergEuler1936, Schwinger1951, Dunne2004}?  The
campaign's design documents call this the Schwinger scaling challenge, and
they gate it behind two contracts.

The first contract is anti-circularity.  A candidate model may not contain
the Schwinger exponent as an input, may not weight or select its initial
conditions by anything equivalent to the target law, may not postselect on
trajectories that reverse, and may not calibrate on the parameter range
used for evaluation.  The second contract is the closure-prescription
clause.  A companion paper of this campaign replicated a published
classical pair-creation event \cite{Land2016} and proved that its reversal
threshold sits exactly on the pole of the Cayley transform that the
published impulsive closure applies to the interaction, while continuous
accumulation of the same interaction provably never reverses
\cite{ZenodoCayley}.  A discretization can manufacture a threshold,
complete with a physically appealing energy condition, that the dynamics
it approximates does not have.  The clause therefore requires continuous
integration through every interaction, a preregistered integration
prescription, agreement between two independent prescriptions at every
claim-bearing point, and stability under refinement.

This paper reports the first sealed passage through that gate.  The
declared family is a resolved Sauter-slab model in which a hidden thermal
oscillator can transfer momentum to the projected time coordinate while
the event crosses a field slab of finite width.  The family was reached
through a four-iteration pilot in which every failure became a design
constraint, and the pilot's central discovery is that the family has a
measured critical manifold in its parameter plane.  Above the manifold a
deterministic crossing kick reverses every trajectory.  Below it reversal
is a rare fluctuation, and a thermal-free probe measures an effective gap
that organizes the suppression.

Two sealed runs then interrogated that structure.  The first asked whether
the effective-gap axis predicts held-out cells to binomial precision and
whether the Schwinger-shaped axis fails by a fixed factor.  Its verdict
was neither-axis.  The effective-gap model organized the data well and
beat the alternative decisively, but its residuals exceeded binomial
scatter by far, so the sealed bars failed.  The bars had tested a claim
nobody made, that a single-axis phenomenological law holds to counting
precision at fifty thousand trajectories per cell.  The discipline caught
its own designer, and the verdict stands in the record.  Before any
second attempt, the project owner imposed a standing rule.  A second
preregistration had to either model the two-variable structure of the
exponent explicitly or set its bars against model error, and it had to
declare which in writing first.  The second sealed run took the second
branch, on grids fully disjoint from everything measured before, and
passed both bars with room.

The result is a disciplined negative.  In this family the reversal
exponent is set by the measure of the initial ensemble, which is exactly
the outcome the campaign's generic negative control anticipated, and the
Schwinger-shaped axis fails to compete.  Section \ref{sec:family} defines
the family and its measured structure.  Sections \ref{sec:run1} and
\ref{sec:run2} report the two governed runs.  Section \ref{sec:meaning}
states what is and is not established, separated by claim type.  Section
\ref{sec:methods} records the evidence infrastructure.

\section{The family and its measured structure}
\label{sec:family}

\subsection{Definition}

The Sauter-slab energy-transfer family evolves an event coordinate $t$
and a hidden oscillator coordinate $u$ under the Hamiltonian
\begin{equation}
H = \tfrac{1}{2}p_{t}^{2} + \tfrac{1}{2}p_{u}^{2}
  + \tfrac{1}{2}\omega_{u}^{2}u^{2} + \tfrac{\lambda}{4}u^{4}
  + g E L \tanh(t/L)\, u,
\label{eq:ham}
\end{equation}
with $\omega_{t} = 0$, so the time coordinate is free except for the
coupling.  The force on the time momentum is
$-gEu\,\mathrm{sech}^{2}(t/L)$, a field localized in a slab of width $L$,
so the field does bounded work on the event.  The committed parameters
are $(\omega_{u}, \lambda, g, T, L) = (1.2,\, 0.1,\, 0.25,\, 1,\, 3)$.
The event enters from $t(0) = -4L$ with momentum gap $p_{t}(0) = P > 0$.
The hidden oscillator is initialized thermally at temperature $T$ about
the Newton-solved shifted equilibrium of its effective well under the
saturated tilt.  Integration is velocity Verlet at $dt = 10^{-3}$ for forty units of
internal time, and the scheme is symplectic here because every force
depends on position only \cite{HairerLubichWanner}.  Reversal means
$p_{t} \le 0$ at any step.
The family contains no exponential in its inputs, no field-dependent
sampling, and no stopping rule correlated with reversal, in compliance
with the anti-circularity contract.

\subsection{Four iterations, each failure a constraint}

The pilot reached this definition in four steps, and the record preserves
all of them.

First, the original coupling $gEtu$ is unbounded in $t$, so secular
growth let the field do unlimited work and the relative energy drift
reached $0.23$ at the strongest field.  The slab replaced it, which is
why the family conserves energy to $6\times10^{-7}$ everywhere below.

Second, the slab version initialized the oscillator thermally about
$u = 0$ while the saturated tilt already displaces the equilibrium at
$t(0)$.  Every ensemble member then carried the same coherent transient,
whose deterministic momentum transfer reversed every trajectory, a
measured fraction of $1.0$.  Nothing was a rare event.  The fix
initializes the ensemble about the shifted equilibrium.

Third, with the transient removed, a deterministic momentum transfer
remained during the slab crossing, and it feeds back on itself, since a
slowing event spends longer in the slab and absorbs more kick.  The
family therefore has a critical-field manifold $E_{\mathrm{crit}}(P)$
separating deterministic reversal from fluctuation-driven reversal.  The
pilot turned this from a nuisance into an instrument.  A thermal-free
probe trajectory, started at the shifted equilibrium with no oscillator
momentum, measures the deterministic minimum $d(P, E)$ of the time
momentum in each cell and classifies its regime.  Cells with $d \le 0$
reverse deterministically and are excluded from every rare-event fit.
For $d > 0$ the probe minimum is the family's effective gap, the
momentum the fluctuations must supply.

Fourth, the sub-critical transition is razor sharp.  At gap $P = 0.5$
the measured fraction moved from $2.8\times10^{-4}$ to $0.91$ between
fields $0.65$ and $0.80$, three decades of rate in a fifteen percent
field step, and only two of the eighteen pilot cells landed in the
measurable band.  Fixed rectangular grids mostly buy bounds in this
family.  Both governed runs therefore place their cells adaptively, by
bisecting the field until the probe hits declared targets in $d/E$,
which uses no ensemble and no rate information.

\begin{figure}[t]
\centering
\begin{tikzpicture}
\begin{axis}[paperaxis, height=0.62\columnwidth,
  xlabel={field $E$},
  ylabel={probe minimum $d$},
  xmin=0.25, xmax=1.05, ymin=-0.6, ymax=1.0,
  legend pos=south west]
\fill[red!8] (axis cs:0.25,-0.6) rectangle (axis cs:1.05,0);
\node[font=\footnotesize, text=red!60!black]
  at (axis cs:0.45,-0.42) {deterministic reversal};
\addplot[blue!70!black, mark=*, mark size=1.6pt,
  discard if not={P}{0.5}]
  table[x=E, y=d] {figdata/pf4_manifold.dat};
\addlegendentry{$P = 0.5$}
\addplot[red!70!black, mark=square*, mark size=1.6pt,
  discard if not={P}{0.7}]
  table[x=E, y=d] {figdata/pf4_manifold.dat};
\addlegendentry{$P = 0.7$}
\addplot[black!60, mark=triangle*, mark size=1.8pt,
  discard if not={P}{0.9}]
  table[x=E, y=d] {figdata/pf4_manifold.dat};
\addlegendentry{$P = 0.9$}
\addplot[only marks, mark=o, mark size=3.2pt, red!70!black,
  discard if not={regime}{1}, forget plot]
  table[x=E, y=d] {figdata/pf4_manifold.dat};
\end{axis}
\end{tikzpicture}
\caption{The measured critical manifold of the pilot.  Each point is the
thermal-free probe's deterministic minimum $d$ of the time momentum for
one pilot cell, plotted against the field for the three pilot gaps.  The
shaded region marks deterministic reversal, where the probe itself
crosses zero and every trajectory reverses without fluctuations.  The
circled point is the one pilot cell in that regime.  For $d > 0$ the
probe minimum is the effective gap that organizes the rare-event
suppression.  Data are the probe entries of the committed pilot record.}
\label{fig:manifold}
\end{figure}

\subsection{Pilot suppression structure}

The two measurable pilot cells admit one diagnostic, stated here with its
caveat attached.  Their suppressions satisfy
$-\log f \propto (d/E)^{2}$ to two percent, a measured ratio of $2.03$
against a predicted $1.99$, while the Schwinger-shaped axis $d^{2}/E$
misses by fifty-seven percent and the bare-gap axis $(P/E)^{2}$ by forty
percent.  Two points prove nothing, and this diagnostic is post hoc and
exploratory.  What it did was point the sealed design at the effective-gap
Gaussian axis, which is the form a measure tail takes when reversal
requires the thermal ensemble to supply momentum $d$ against fluctuations
whose scale is set by $E$, the classical escape-rate structure of Kramers
\cite{Kramers1940}.  The campaign's generic negative control had already
found that fold statistics in a bounded Hamiltonian family are set
rigidly by the dynamics and the initial-condition measure
\cite{ZenodoEntropy}, so a measure-tail exponent was the expected shape.

\section{The first sealed run}
\label{sec:run1}

\subsection{Design}

PREREG-PF4-001 declared the claim that the fluctuation-regime suppression
is organized by the effective-gap Gaussian axis and not by the
Schwinger-shaped axis.  The two models are
\begin{align}
\text{model G}\qquad -\log f &= \alpha + \beta\,(d/E)^{2},
\label{eq:modelg}\\
\text{model S}\qquad -\log f &= \alpha' + \beta'\,d^{2}/E,
\label{eq:models}
\end{align}
fit by weighted least squares with binomial weights
$\mathrm{var}(-\log f) = (1-f)/(fN)$ on twelve training cells and scored
on eight held-out cells.  Training gaps were $P \in \{0.5, 0.7, 0.9\}$
and held-out gaps $P \in \{0.6, 0.8\}$, disjoint in $P$, with four cells
per gap placed by bisection at probe targets
$d/E \in \{0.20, 0.28, 0.36, 0.44\}$ and $5\times10^{4}$ trajectories
per cell.  The sealed bars were $\mathrm{MSE}_{G} \le 4$ in binomial
units and $\mathrm{MSE}_{S} \ge 4\,\mathrm{MSE}_{G}$.  The claim would be
refuted if model S matched or beat model G held out.  Any other outcome
was to be reported as neither-axis with both residual sets.

\subsection{Outcome, verdict preserved}

The instruments worked.  All twenty probe-placed cells landed in the
measurable band, so the adaptive manifest closed the pilot's grid
problem.  Drift was uniform at $6\times10^{-7}$.  Both prescription
controls agreed exactly, with $z = 0.00$ under an independent
Runge-Kutta integration and under step halving.  The fitted
effective-gap slope was $\beta = 41.2$, confirming the pilot's two-point
prediction of $41.6$, and model G beat model S by a factor of $3.28$ on
held-out weighted error.

The sealed bars were still not met.  The verdict was neither-axis,
because $\mathrm{MSE}_{G} = 81.2$ against a bar of $4$, and model S,
though clearly worse, was not four times worse.  At
$5\times10^{4}$-trajectory precision the residual dependence of the
suppression on the gap at fixed $d/E$ is many binomial standard
deviations.  Stated plainly, the sealed bars tested a claim nobody made,
that a single-axis phenomenological law holds to binomial precision.
The effective-gap axis is a few-percent-accurate organizer of this
family and not a law, and bars written in counting units could only
discover that fact by failing.  The discipline caught its own designer.
The verdict is preserved in the record exactly as returned, and nothing
in the second run modifies it.

\section{The correction rule and the second sealed run}
\label{sec:run2}

\subsection{The owner-imposed rule}

Before any second attempt existed, the design document acquired a
standing requirement.  A second sealed attempt on this family must
either model the two-variable structure explicitly, an exponent
depending on more than one combination of $P$ and $E$, or set its bars
against model error rather than binomial error, and the preregistration
must declare which of the two it does before any governed run executes.
The rule was declared in writing before PREREG-PF4-002 was drafted, and
the second preregistration opens by taking the second branch.  Its bars
are ratios to the model's own training inadequacy, so the claim under
test is the one actually made, that the effective-gap axis organizes and
generalizes while the Schwinger-shaped axis does not.

\subsection{Design}

PREREG-PF4-002 kept the family, the estimator, and the control rules of
the first run unchanged and changed only the grids and the bar
arithmetic.  Training gaps were $P \in \{0.55, 0.75, 0.95\}$ and
held-out gaps $P \in \{0.65, 0.85\}$, disjoint from the pilot grid and
from the first governed run.  Probe targets were
$d/E \in \{0.22, 0.30, 0.38, 0.46\}$, disjoint from the first run's
targets.  The bars were that model G's held-out weighted error be at
most twice its training weighted error, and that model S's held-out
weighted error be at least twice model G's.

\subsection{Outcome}

Both bars passed with room.  The generalization ratio was $0.38$, so
model G predicted the gaps it had never seen better than it fit the gaps
it trained on.  The Schwinger-shaped model was $3.14$ times worse on the
same held-out cells.  Ten of twelve training cells and seven of eight
held-out cells were usable under the sealed count and saturation rules,
above the declared minimum.  Drift was uniform at $6\times10^{-7}$.
The verdict is a pass, and the claim carries the label
demonstrated-in-model.

Both prescription controls again agreed exactly, with $z = 0.00$ under
Runge-Kutta and under step halving, as they had in the first run.  This
exactness is an observation in its own right.  The control cells sit at
fractions near $0.01$ and $0.1$, where identical fractions across three
prescriptions mean the integrators agree member for member on which of
$5\times10^{4}$ trajectories reverse, not merely on the ensemble rate.
At this step size the family's reversal classification is robust per
trajectory, which is the strongest form of the closure clause's
agreement requirement and the opposite of the prescription-borne
threshold the clause exists to catch.

\begin{figure}[t]
\centering
\begin{tikzpicture}
\begin{axis}[paperaxis, height=0.62\columnwidth,
  xlabel={$(d/E)^{2}$},
  ylabel={$-\log f$},
  legend pos=north west]
\addplot[blue!70!black, only marks, mark=*, mark size=1.6pt,
  discard if not={split}{0}]
  table[x=xg, y=neglogf] {figdata/pf4_cells.dat};
\addlegendentry{training cells}
\addplot[red!70!black, only marks, mark=square*, mark size=1.7pt,
  discard if not={split}{1}]
  table[x=xg, y=neglogf] {figdata/pf4_cells.dat};
\addlegendentry{held-out cells}
\addplot[black, thick]
  table[x=xg, y=pred] {figdata/pf4_gline.dat};
\addlegendentry{sealed G model, run 002}
\end{axis}
\end{tikzpicture}
\caption{Suppression against the effective-gap Gaussian axis for every
usable cell of both governed runs, training and held-out marked
separately, with the second run's sealed G-model line overlaid.  The
axis organizes thirty-six cells spanning more than three decades of
rate across six trained and four held-out gap values.  The visible scatter about
the line is real structure, many binomial standard deviations at this
ensemble size, which is why the first run's binomial-unit bars failed
and the second run's bars are ratios to model error.  Data are the cell
entries of the two committed run records.}
\label{fig:cells}
\end{figure}

\section{What the result means and does not mean}
\label{sec:meaning}

\subsection{Claim register}

The campaign separates claim types, and this paper's claims sort as
follows.

Theorem-level.  The fold kinematics behind apparent pair events, the
exact conservation argument that makes the static-potential family an
analytic null with no dynamical exponential, and the frozen-kick result
that impulsively closed interactions cannot produce continuous reversal
\cite{ZenodoCayley} are proved statements, not measurements.

Sealed and measured.  The neither-axis verdict of PREREG-PF4-001 and the
passing verdict of PREREG-PF4-002 are governed outcomes under sealed
bars, with the second carrying the label demonstrated-in-model.  Together
they establish that in the declared Sauter-slab family the reversal
suppression is organized by the effective-gap Gaussian axis within its
own accuracy class, generalizing across gaps it never trained on, and
that the Schwinger-shaped axis fails to compete.

Exploratory.  The pilot sweep, the critical-manifold map, the two-point
diagnostic, and the slope confirmation $\beta = 41.2$ against the
predicted $41.6$ are unsealed or secondary observations.  They shaped
the design and none of them bears the claim.

Open.  Whether any constructible hidden-dynamics family has a reversal
exponent that is not a measure tail is not addressed by any run reported
here.

\subsection{Meaning}

The exponent in this family is measure-set.  Reversal happens when the
thermal ensemble supplies the effective gap, so the suppression is the
Gaussian tail of the declared initial measure, evaluated in a variable
the dynamics itself defines through the probe.  This is exactly the
shape the campaign's generic negative control anticipated for bounded
local Hamiltonian families, now established under seal rather than
observed in passing.  The Schwinger form, which differs from the
measure-tail form by one power of the field in the exponent, had a fair
preregistered chance on held-out data twice and failed to compete both
times.

\subsection{Limits}

The verdict binds one declared family at its declared parameters.  It
licenses no statement about other hidden-dynamics families, none about
physical pair production, and none about quantum electrodynamics, whose
exponent is derived structure, not a fitted axis
\cite{Schwinger1951, Dunne2004}.  The open summit of the program is now
sharply posed.  Either some constructible family has a reversal exponent
that is not the tail of its own initial measure, or the Schwinger
exponential is unreachable by this route.  Any claimed positive would
owe its audience what this paper's negative paid, continuous integration
under the closure clause, prescription-pair and refinement controls at
every fitted point, the planted-artifact check, disjoint held-out grids,
and bars declared against the claim actually made.

\section{Methods and reproducibility}
\label{sec:methods}

The campaign's instruments were sealed under tolerance freezes before
any claim-bearing run, and the seal ledger now carries six entries, two
instrument freezes and four preregistrations, each a triple of
registration identifier, sealing commit, and content hash.  Evidence
records are append-only JSON with SHA-256 hashes over their own
content, and the three records behind this paper are
\texttt{results/pf4-pilot.json},
\texttt{results/prereg-pf4-001.json}, and
\texttt{results/prereg-pf4-002.json}.
Every number quoted here appears verbatim in one of those records or in
the sealed design documents, and both figures are generated by a
committed script whose only inputs are the committed records.

The closure clause was active throughout.  All dynamics were integrated
continuously with velocity Verlet at the declared step, the prescription
was preregistered, and the prescription-pair and refinement controls ran
on cells verified measurable in advance, closing a pilot lesson in which
a control cell with zero counts had certified nothing.  The planted
Cayley-trap control, clause C5, was deferred in both preregistrations
with the recorded justification that these runs compare axes within one
declared continuous prescription rather than asserting a rate value
against an external target, while C3 and C4 remained active.  The
deferral is not a waiver.  Any future preregistration that claims a rate
law against a physical target must implement the full trap, and a
pipeline that fails to flag the planted artifact voids its own run.

\begin{thebibliography}{10}

\bibitem{Sauter1931}
F. Sauter,
\"Uber das Verhalten eines Elektrons im homogenen elektrischen Feld
nach der relativistischen Theorie Diracs,
Z. Phys. \textbf{69}, 742 (1931).

\bibitem{HeisenbergEuler1936}
W. Heisenberg and H. Euler,
Folgerungen aus der Diracschen Theorie des Positrons,
Z. Phys. \textbf{98}, 714 (1936).

\bibitem{Schwinger1951}
J. Schwinger,
On gauge invariance and vacuum polarization,
Phys. Rev. \textbf{82}, 664 (1951).

\bibitem{Dunne2004}
G.~V. Dunne,
Heisenberg-Euler effective Lagrangians: Basics and extensions,
in \textit{From Fields to Strings: Circumnavigating Theoretical
Physics}, Ian Kogan Memorial Collection
(World Scientific, Singapore, 2005),
arXiv:hep-th/0406216.

\bibitem{Kramers1940}
H.~A. Kramers,
Brownian motion in a field of force and the diffusion model of
chemical reactions,
Physica \textbf{7}, 284 (1940).

\bibitem{Land2016}
M. Land,
Pair production in classical Stueckelberg--Horwitz--Piron
electrodynamics,
J. Phys.: Conf. Ser. \textbf{615}, 012007 (2015),
arXiv:1604.01625.

\bibitem{HairerLubichWanner}
E. Hairer, C. Lubich, and G. Wanner,
\textit{Geometric Numerical Integration}, 2nd ed.
(Springer, Berlin, 2006).

\bibitem{ZenodoEntropy}
A.~H. Bond,
Entropy across singular projections: Fiber multiplicity, caustics, and
observer-relative information,
Zenodo (2026), doi:10.5281/zenodo.21789011.

\bibitem{ZenodoCayley}
A.~H. Bond,
The pair-creation threshold of impulsive Stueckelberg--Horwitz--Piron
scattering is a Cayley-transform pole,
Zenodo (2026), doi:10.5281/zenodo.21790095 (v2).

\bibitem{ZenodoQO}
A.~H. Bond,
Consumer-relative quantum distinguishability on causal boundaries,
Zenodo (2026), doi:10.5281/zenodo.21789046.

\end{thebibliography}

\end{document}
